实际扫描取301个等距$B$值（300区间）, RK4每周期256步, 初值$(.2,0)$, 先舍弃500周期再观察100周期, 图取后32周期. 选例另外用自适应积分复算, 相对容差$2\times10^{-10}$、绝对容差$2\times10^{-12}$、最大步长$T_D/64$, 舍弃900周期后观察120周期. 得到下表, 回归误差是最后40对截面点的最大状态距离.
\begin{center}
\begin{tabular}{rlll}
\toprule $B$ & 观察周期 & 回归误差 & 有限时间指数\\\midrule
9.6 & 2 & $3.60\times10^{-10}$ & $-.04992$\\
9.8666667 & 4 & $1.18\times10^{-11}$ & $-.03895$\\
7.9333333 & 未检出$\le16$周期 & 不小于$3.69$ & $.07240$\\
10.1833333 & 未检出$\le16$周期 & 不小于$4.03$ & $.06919$\\\bottomrule
\end{tabular}
\end{center}
前两例分别在二、四次截面迭代后回归, 且更低周期回归误差不小, 是稳定倍周期与四倍周期的强数值证据. 后两例本次点集未收敛为有限个重复点、切向增长率为正; 第三例两近初值在$t\le659.73$内状态距离最大达到8.13. 这些结果支持原题所需的两个不同混沌数值案例, 并清楚保留有限精度与有限时观察的边界. 自适应复算与RK4扫描使用不同积分方法, 对混沌轨道不要求逐点长期重合.
